\begin{table}[H]
\centering
\caption{Tasas laborales por logro educativo, 2010 y 2025}
\label{tab:tasas-laborales-educacion}
\begin{threeparttable}
\scriptsize
\setlength{\tabcolsep}{3.2pt}
\resizebox{\textwidth}{!}{%
\begin{tabular}{lrrrrrrrr}
\toprule
& \multicolumn{2}{c}{Participación}
& \multicolumn{2}{c}{Ocupación}
& \multicolumn{2}{c}{Desempleo}
& \multicolumn{2}{c}{Formalidad} \\
\cmidrule(lr){2-3} \cmidrule(lr){4-5} \cmidrule(lr){6-7} \cmidrule(l){8-9}
Logro educativo & 2010 & 2025 & 2010 & 2025 & 2010 & 2025 & 2010 & 2025 \\
\midrule
Ninguno o preescolar & 51,5 & 40,8 & 47,8 & 39,0 & 7,1 & 4,3 & 4,1 & 3,9 \\
Básica primaria & 61,8 & 55,4 & 56,6 & 52,4 & 8,3 & 5,5 & 11,7 & 13,9 \\
Básica secundaria & 45,5 & 55,1 & 39,7 & 50,3 & 12,6 & 8,7 & 17,6 & 18,3 \\
Educación media & 73,5 & 67,6 & 62,6 & 60,8 & 14,8 & 10,0 & 36,3 & 41,2 \\
Superior o universitaria & 78,5 & 74,3 & 68,4 & 66,9 & 12,9 & 10,0 & 65,4 & 75,4 \\
\midrule
\textbf{Total} & 62,9 & 64,3 & 55,6 & 58,6 & 11,7 & 8,9 & 30,4 & 44,0 \\
\bottomrule
\end{tabular}
}
\begin{tablenotes}
\scriptsize
\item \textit{Nota:} La tasa de participación se calcula como PEA/PET, la tasa de ocupación como ocupados/PET, la tasa de desempleo como desocupados/PEA y la tasa de formalidad como ocupados que cotizan a pensión sobre ocupados. Las tasas se reportan en porcentaje. Fuente: cálculos propios con datos GEIH del DANE.
\end{tablenotes}
\end{threeparttable}
\end{table}
